\topic{本章实验、故意破坏与练习}

\begin{experimentbox}{逐 byte 审计复位向量}
构建 ROM 后，用逐 byte 工具查看偏移 \texttt{0x1FFF0}，记录文件大小、起始
偏移和 \texttt{E9 0D F0}。纸上把 rel16 解释为补码，再算目标 IP。
\end{experimentbox}

\begin{faultinjectionbox}{把位移 byte 交换}
把 \texttt{0D F0} 当作 \texttt{F0 0D}。预期目标不再是
\texttt{F000}，最早串口日志消失。恢复 little-endian 解释后重新计算。
\end{faultinjectionbox}

\begin{pitfallbox}
容量不是最后地址；byte 顺序不是 bit 顺序；补码值必须先固定宽度。
\end{pitfallbox}

\textbf{练习}
\begin{enumerate}[leftmargin=2.2em]
  \item 12 bit 能编码多少状态，最大无符号值是多少？
  \item \texttt{FE} 的 8-bit 补码值是什么？
  \item \texttt{78 56} 按 little-endian 读成什么？
\end{enumerate}
\begin{answerbox}
1. \(4096\) 种、最大 4095；2. \(254-256=-2\)；3. \texttt{0x5678}。
\end{answerbox}
